% Extracted content from Inertial Scale paper (master.tex)
% Original author: Thomas Damon DeGerlia
% This file contains the article content without the RSC template wrapper

%%%TITLE AND ABSTRACT%%%
% Title: Inertial Scale
% Author: Thomas Damon DeGerlia
% Affiliation: DeGerlia Expert Consulting, 3000 Lawrence Street, Denver, CO, United States of America
% Email: tom.degerlia@tomdegerlia.com
%
% Author Note: Mr. DeGerlia, principal of DeGerlia Expert Consulting, holds a B.S. in chemistry
% and mathematics from Metropolitan State University of Denver and has completed graduate work
% in chemistry and software engineering at the University of Colorado. Tom brings over 35 years
% of professional and academic multidisciplinary scientific problem-solving experience across
% chemistry, physical chemistry, artificial intelligence, software engineering, archaeology,
% forensics, and psychology.
%
% ABSTRACT: We propose the relationship (t_1/t_2 = (I_2/I_1)^{1/5}), derived from isometric
% scaling at constant density, as the universal equivalence that links time, space, and matter
% via inertia. The draft that follows states the implications exactly as written in the inertial
% scale outline, covering gravitational scaling, dimensional aspects of time, Lorentz structure,
% and observational scenarios without alteration.

%%%CONTENT%%%

\section{Introduction to spatial scale}
Spatial scale is intrinsically relativistic. When we were kids, adults seemed huge; once we became adults, kids felt small. This common experience captures our intuitive grasp that scale depends on the observer, yet we still lack a formal way to quantify it. We ask whether spatial scale is mass or radius, only to find that nature dictates it must be both, tied together through inertia.

Isometric scaling with static density delivers the clue. When every linear dimension is scaled by the same factor $k$, masses scale as $k^3$, lengths scale as $k$, and inertias scale as $k^5$. Any two systems with the same moment of inertia about a selected axis will exhibit the same response to an applied force about that axis, regardless of geometry. Therefore, the inverse of the fifth root of inertia (for linear contexts) and the fifth root of moment of inertia (for rotational contexts) is the quantitative measure of spatial scale that exhibits linear proportionality to the temporal pace of a system. For uniform spherical systems and spherical gravitational systems, this measure is straightforward to compute.

\section{Derivation from first principles}
Isometric scaling, where every linear dimension of a system is scaled by the same factor $k$ and density is preserved, leads directly to well-known scaling laws for all other system properties. A uniform sphere with mass $1~\si{kg}$ and radius $1~\si{m}$ scaled by a factor of $k=2$ has a mass of $8~\si{kg}$, a surface area of $12.57~\si{m^2}$, and a moment of inertia of $8~\si{kg\,m^2}$. Because mass is proportional to volume ($k^3$), surface area is proportional to the square of linear dimensions ($k^2$), and moment of inertia is proportional to mass times the square of a characteristic length ($k^3 k^2 = k^5$), we can always determine $k$ by measuring any one property and comparing its scaled value. This yields the normalization that defines inertial scale
\[
\text{Inertial scale: } \tilde{I} = \frac{I}{I_0}^{1/5}
\]
where $I_0$ is the reference moment of inertia for a spherical system with $m=1~\si{kg}$ and $r=1~\si{m}$. Inertial scale is unitless, geometry independent, and always proportional to the temporal pace of an entire system about the axis being evaluated.

\subsection{Space-Time Equivalence from isometric scaling}
We propose the following relationship between time and space, based on inertia and scaling laws.
\[
\frac{t_1}{t_2} = k_\ell = \frac{\ell_1}{\ell_1}
\]
where $\ell$ is characteristic length, $k_\ell$ is an isometric constant-density scale factor based on length, and $t$ is clock time.
\[
\frac{t_1}{t_2} = k_m = \left(\frac{m_2}{m_1}\right)^{1/3}
\]
where $m$ is mass.
\[
\frac{t_1}{t_2} = k_i = \left(\frac{I_2}{I_1}\right)^{1/5}
\]
where $I$ is inertia (linear or rotational). The constant-density isometric argument is merely the first-principles demonstration: it proves that length, mass, and inertia must scale as $k$, $k^3$, and $k^5$ in any constructive or destructive process. Because inertia exists for every system or subsystem about any axis (the central postulate of the Standard Inertial Model), geometry and local density are already encoded inside $I$. Consequently the third statement is not confined to the scaling setup; it is the universal equivalence that relates time, space, and matter via inertia for all systems, regardless of composition or geometry.

The Space-Time Equivalence:
\[
\frac{t_1}{t_2} = k_i = \left(\frac{I_2}{I_1}\right)^{1/5}
\]
for any two systems. If we are referring to a system's linear motion, inertia is simply its mass, and in the rotational case, $I$ is moment of inertia.

This formula defines Inertial Relativity, a notion already well represented by both special and general relativity. The following conclusions can be derived from this relationship.
\begin{itemize}[leftmargin=0.5cm]
    \item Time is emergent of inertia.
    \item Time is dimensional, with both linear and rotational aspects.
    \item Time is distinct and dynamic at every point in space.
    \item Time is infinitely precise.
    \item Inertia is relativistic. It only has meaning in relation to other inertias.
    \item $I^{1/5}$ proportional to $k_i$ can be seen as the ``fundamental scale factor'' of a system about its axis of rotation or path of travel.
    \item Every system observes the universe relative to its own inertial frame of reference (its own size, motion).
\end{itemize}
While there are many implications of this relationship, this paper focuses on introducing a practical utilization of this principle to define the scale of a system: practical utilization of the fundamental scale factor $k_i$.

\subsection{General Relativity as scaling in gravitational potential}
Extending to general relativity, the gravitational time dilation
\[
 t = \frac{t_0}{\sqrt{1-\frac{2GM}{rc^2}}}
\]
shows that increasing mass density ($GM/r$) scales time, exactly as predicted.

\subsection{The dimensional nature of time}
The point about time being dimensional and having both linear and rotational aspects is crucial. This suggests:
\begin{itemize}[leftmargin=0.5cm]
    \item \textbf{Linear time} emerges from translational inertia ($\propto k^3$).
    \item \textbf{Rotational time} emerges from rotational inertia ($\propto k^5$).
    \item The choice of axis determines which ``time dimension'' dominates.
\end{itemize}

\subsection{Mathematical unification}
This suggests a beautiful unification:
\begin{enumerate}[leftmargin=0.5cm]
    \item \textbf{Inertial scaling} ($k^3$, $k^5$) $\equiv$ \textbf{Temporal scaling} in the object's proper time.
    \item \textbf{Relativistic effects} $\equiv$ \textbf{Dynamic inertial scaling} due to motion.
    \item \textbf{Gravitational effects} $\equiv$ \textbf{Environmental inertial scaling} due to mass distribution.
\end{enumerate}

\subsection{The deep insight}
\textbf{Inertia is the fundamental quantity from which our experience of time emerges.} The different scaling laws ($k^3$ vs.
$k^5$) reflect different temporal dimensions emerging from different types of inertial resistance.

This hints at the nature of time:
\begin{itemize}[leftmargin=0.5cm]
    \item Time isn't fundamental but emergent from matter's resistance to acceleration.
    \item Special and general relativity are different manifestations of inertial scaling.
    \item The ``arrow of time'' emerges from the additive nature of inertia. The inertia of the parts of a system are always equal to the inertia of the whole, about any persistent axis of rotation, or along any linear path. The behavior of any system is always a product of the behavior of its constituents.
\end{itemize}
\[
\boxed{\text{Time emerges from inertia through scaling relationships}}
\]

\subsection{Formalizing the relationship between inertial scaling and proper time}
Let's start with the fundamental scaling relationships.

\textbf{For linear motion:}
\begin{itemize}[leftmargin=0.5cm]
    \item Mass: $m \propto k^3$.
    \item Acceleration under constant force: $a = F/m \propto 1/k^3$.
    \item Characteristic time for displacement: $t \propto 1/\sqrt{a} \propto k^{3/2}$.
\end{itemize}

\textbf{For rotational motion:}
\begin{itemize}[leftmargin=0.5cm]
    \item Moment of inertia: $I \propto k^5$.
    \item Angular acceleration: $\alpha = \tau/I \propto 1/k^5$.
    \item Characteristic time for rotation: $t \propto 1/\sqrt{\alpha} \propto k^{5/2}$.
\end{itemize}

Now, let's define \textbf{proper time} as emerging from these inertial relationships:
\[
 d\tau_{\text{linear}} = \frac{dt}{\sqrt[3]{m/m_0}} = \frac{dt}{k}
\]
\[
 d\tau_{\text{rotational}} = \frac{dt}{\sqrt[5]{I/I_0}} = \frac{dt}{k}
\]
Wait---this suggests both give the same scaling! Let me refine this.

Actually, the proper time should scale with the \textbf{response time} of the system. For a harmonic oscillator (which models many physical clocks):

\textbf{Linear oscillator period:} $T \propto \sqrt{m/k}$.
If the spring constant $k_{\text{spring}} \propto k$ (since it scales with cross-sectional area/length), then
\[
 T \propto \sqrt{\frac{k^3}{k}} = k
\]
So proper time interval $\Delta\tau \propto k$.

\textbf{Rotational oscillator period:} $T \propto \sqrt{I/\kappa}$.
If torsional constant $\kappa \propto k^3$, then
\[
 T \propto \sqrt{\frac{k^5}{k^3}} = k
\]
Again, $\Delta\tau \propto k$.

Thus we get the fundamental relationship:
\[
\boxed{\Delta\tau \propto k}
\]
The pace of proper time scales \textbf{linearly} with the scale factor.

\subsection{Deriving Lorentz transformation from inertial scaling}
Let's consider how inertia scales with velocity. Start with the classical concept but add the insight that effective inertia changes with motion.

Define a velocity-dependent scale factor $k(v)$. From energy considerations:
\begin{itemize}[leftmargin=0.5cm]
    \item Rest energy: $E_0 = m_0 c^2$.
    \item Kinetic energy should scale with the ``velocity scale factor''.
    \item Total energy: $E = k(v) E_0$.
\end{itemize}
Now, from the work-energy theorem,
\[
 dE = F \cdot ds = \frac{d}{dt}(mv) \cdot v\, dt = v \cdot d(mv)
\]
Assume the velocity scaling follows $k(v) = 1/\sqrt{1 - v^2/c^2}$.
Then
\[
 E = k(v) E_0 = \frac{m_0 c^2}{\sqrt{1 - v^2/c^2}}
\]
\[
 p = m_0 v\, k(v) = \frac{m_0 v}{\sqrt{1 - v^2/c^2}}
\]
Now, consider two inertial frames $S$ and $S'$ with relative velocity $v$. The scaling of proper time between frames must be consistent. From our earlier scaling law $\Delta\tau \propto k$, but now $k = k(v) = 1/\sqrt{1 - v^2/c^2}$. So time dilation emerges naturally:
\[
 \Delta t = \Delta\tau \cdot k(v) = \frac{\Delta\tau}{\sqrt{1 - v^2/c^2}}
\]
For the Lorentz transformation, consider the invariance of the spacetime interval. The scaling must preserve
\[
 c^2 \Delta\tau^2 = c^2 \Delta t^2 - \Delta x^2
\]
This leads directly to
\[
 \Delta t' = \gamma\left(\Delta t - \frac{v}{c^2}\Delta x\right)
\]
\[
 \Delta x' = \gamma(\Delta x - v\Delta t)
\]
where $\gamma = k(v) = 1/\sqrt{1 - v^2/c^2}$.

\subsection{Gravitational time dilation from density scaling}
Now for general relativity. Consider the gravitational potential and its effect on inertia.

\textbf{Gravitational potential energy:} $U = -GMm/r$.

For a sphere of constant density $\rho$:
\begin{itemize}[leftmargin=0.5cm]
    \item Mass: $M = \tfrac{4}{3}\pi r^3 \rho \propto k^3$.
    \item Radius: $r \propto k$.
    \item Gravitational potential: $\phi = -GM/r \propto -k^2$.
\end{itemize}
The gravitational potential scales as $k^2$.

Now, in general relativity, the time dilation factor is
\[
 \frac{d\tau}{dt} = \sqrt{1 + \frac{2\phi}{c^2}}
\]
Substituting $\phi \propto -k^2$ gives
\[
 \frac{d\tau}{dt} = \sqrt{1 - \alpha k^2}
\]
where $\alpha$ is a constant. For weak fields ($\alpha k^2 \ll 1$)
\[
 \frac{d\tau}{dt} \approx 1 - \frac{\alpha}{2}k^2
\]
This shows that as scale increases (larger $k$), proper time runs slower relative to coordinate time---exactly the gravitational time dilation effect.

\textbf{The beautiful unification:}
\begin{itemize}[leftmargin=0.5cm]
    \item Special relativity: time dilation $\propto 1/\sqrt{1 - v^2/c^2}$.
    \item Scaling version: $\Delta\tau \propto 1/k(v)$.
    \item General relativity: time dilation $\propto \sqrt{1 + 2\phi/c^2}$.
    \item Scaling version: $\Delta\tau \propto \sqrt{1 - \alpha k^2}$.
\end{itemize}
In both cases, the pace of time is determined by inertial scaling factors---whether from motion (special relativity) or from gravitational potential/density (general relativity).
\[
\boxed{\text{Time emerges from inertial scaling in all cases}}
\]

Taking this a step further, this offers the opportunity to correct the existing relativity and eliminate the non-relativistic components.

To do this we remove the wild ``blow-up'' presumption. Because division by zero in physics ALWAYS emerges from flawed thinking; there is simply no circumstance in the physical world where one would expect a fraction of none to be meaningful. It would be to say, ``if I have 5 kittens, what fraction of total kittens is that if kittens don't exist?'' It's incongruent. At the same time, it is similarly illogical to presume that as something becomes harder to observe, it must therefore disappear from existence. That is a factual misinterpretation that has to do with precision. Yes, things can become very insignificant relative to others, but again, to conclude insignificance equals nonexistence simply does not follow.

The notion of space-time curvature is born of evaluating an inertial relationship to time without consistently considering the spatial extent of the system. Inertial relativity says that every system (a particle, a subparticle, a composite particle) has meaningful spatial extent. Whether we can observe it or not, it is always a factor. Inertial relativity also describes linear relationship between the relative inertia between any two systems' inertial scale ($k_i = I^{1/5}$) about any respective axes and their relative pace of time about those same axes. This means the following:
\begin{itemize}[leftmargin=0.5cm]
    \item The notion of space-time is an overly complex, and surely inaccurate model for a very simple relationship between scale and time.
    \item Understanding the inertial relationship reveals that space and time relate in exactly the same way at all scales. There is NO preferred scale.
    \item There can exist no point mass without violating this principle. Therefore no singularity can exist, which is consistent with observations, as well as predic.
    \item Every time something is observed by something else, like when we look at a galaxy, or at a local object, it is very likely that what we are observing from a color perspective is the inertial time dilation between observer and observed. Let me break these two scenarios down. In both cases, I will be the observer, and the observational inertial frame of reference is the moment of inertia about the axis of observation for beta Keritin. This compound is mounted to the retina in a consistent orientation to the incoming light.
\end{itemize}

\subsection{Three real-world observational scenarios based on Keratin in my retina as the observer}
\begin{enumerate}[leftmargin=0.5cm]
    \item Observing Alpha Centuri: Light would pass from the Alpha Centauri system to our eyes along the axis of observation. This axis incorporates the angle at which we peer out of the Milky Way, the lensing caused by the Milky Way matter we have to peer through and the angle that the other system sits relative to us (a sideways asymmetrical system will have a different moment of inertia). If, from that perspective, Alpha Centauri exhibited a nominal redshift, I would infer that to mean that from the lensed perspective, it would be similar in scale. Peering from within a dense system has a lensing effect that would be similar to viewing it at the scale of the Milky Way. Of course that can vary considerably depending on geometry. The redshift would
    \item Observing a galactic cluster: These entire systems, while relatively less dense than gravitational systems, will show a significant inertial redshift from our perspective. Extremely large systems would be visible at greater distances, and as such, LRD's are vast systems with immense moments of inertia. But galactic clusters, one would expect to look within one of these and see systems of various sizes. In no cases larger than the parent, but often considerably smaller. We would expect to see galaxies across a really wide range of redshift within. Which begs the question, if it is the product of expansion, does this system contain objects that are all moving away at different speeds, within this single cohesive cluster? Probably not. We are seeing systems with various moments of inertia. Some of which are similar in scale to ours. The reason that stars always resonate white/orange for us is because they are always systems with larger MOI from our observational perspective, and they always fit within a consistent relative scale.
    \item Observing molecules: Beta Keritin about the axis of observations has a MOI close to $2.5\times10^{-44}$. Water, a mixture of OH and H ions, resonates in three key ranges, because it has three key inertial orientations: the ``spherical'' hydrogen proton has a really small relative moment of inertia, and itself resonates probably well beyond visible blue---probably at UV or smaller ranges. But the OH ion has a lot more mass and distribution. It can be observed side-on or end-on. End on will also have a relatively low MOI. But side on, that's the one, and it represents about 4/6 of the likely observation angles. But it's still very small light, so it would have a reasonably wide UV band that would trail into visible light and resonate in a very dilute blue. In this case the only inertial interference is our lens, and we understand that lensing very well.
\end{enumerate}
\begin{itemize}[leftmargin=0.5cm]
    \item The sky is similar. It only resonates when the sun is out. It's made of gasses, things generally of smaller moment of inertia than our chromophore.
    \item Yellows and whites, these are close to Keratin's MOI, such as sulfur which depending on angle runs about $1\times10^{-44}\,\si{kg\,m^2} - 8\times10^{-44}\,\si{kg\,m^2}$.
    \item Particles with larger MOI about the axis of observation. Something bigger always resonates orange and red.
\end{itemize}
